\chapter[Structural DSGE Dynamics]{Structural Deterministic Dynamics in DSGE Filtering}
\label{ch:bf-structural-deterministic-dynamics}

\begin{mdframed}
\textbf{Production warning.}
Do not implement DSGE nonlinear filtering by silently treating every declared
state coordinate as an exchangeable noisy Gaussian coordinate.  Many DSGE
models contain a structural split between shock-driven exogenous states and
endogenous or accounting states that are deterministic conditional on lagged
states, controls, and current shocks.  Collapsing that structure into a generic
additive transition can be harmless for a linear Kalman filter, but it is a
structural design error for sigma-point and particle filters unless an explicit
approximation argument is written down and tested.
\end{mdframed}

\section{The Structural Split}
\label{sec:bf-dsge-structural-split}

The state vector in a DSGE likelihood is not merely a numerical vector.  It
usually carries economic timing.  In the language of the structural state
partition from Chapter~\ref{ch:bf-state-space-contracts}, a common DSGE
representation separates a declared stochastic block from a deterministic-
completion block.  Writing the stochastic block as exogenous state variables
$m_t$ and the deterministic-completion block as endogenous predetermined or
accounting state variables $k_t$, one common representation is:
\begin{align}
  m_t &= T_m(m_{t-1}, \varepsilon_t; \theta), \label{eq:bf-m-transition}\\
  k_t &= T_k(k_{t-1}, m_{t-1}, m_t; \theta), \label{eq:bf-k-transition}\\
  x_t &= (m_t, k_t).
\end{align}
In a perturbation solution, this structure may also be represented by a compact
law of motion
\[
  x_t = h_x x_{t-1} + \eta \varepsilon_t
\]
at first order, with zero or rank-deficient innovation rows for the
deterministic-completion block.  The compact representation is useful, but it
must not erase the timing contract or the structural metadata described in
Chapter~\ref{ch:bf-api-design}.  For mixed structural models, the one-step
predictive law is the pushforward of the joint law of the lagged state and the
current stochastic input through the structural transition map:
\[
  p_\theta(x_t\mid y_{1:t-1})
  =
  \iint \delta\!\left(x_t - F_\theta(x_{t-1},\varepsilon_t)\right)
  p(\varepsilon_t)\,p_\theta(x_{t-1}\mid y_{1:t-1})
  \, d\varepsilon_t \, dx_{t-1}.
\]
In nonlinear filtering, the correct random input is therefore the innovation or
declared stochastic block.  The endogenous block is then generated by the
model-implied deterministic-completion map.  Here ``deterministic'' means no
independent innovation in that coordinate once $(x_{t-1},\varepsilon_t)$ is
fixed; it does not mean zero predictive variance conditional only on the lagged
filtering state.

This is the distinction emphasized in the DSGE filtering literature.  It is
also an instance of the generic BayesFilter structural-transition contract from
Chapter~\ref{ch:bf-state-space-contracts}: the backend integrates over the
declared stochastic block and completes the deterministic block pointwise.
Standard Bayesian DSGE expositions write the likelihood in state-space form
after the model solution has supplied a transition law for predetermined
variables and an observation equation for measured macroeconomic series
\citep{herbst2015bayesian,an2007bayesian}.  Pruned perturbation methods make
the same point in a different language: the first-order component is propagated
stochastically, while the higher-order correction is propagated by a separate
deterministic recursion driven by the lower-order state
\citep{kim2008calculating,andreasen2018pruned}.  The filtering backend must
honor the state-space law supplied by the structural solution, not merely the
dimension of the state vector.

\section{Why Linear Kalman Filtering Can Hide the Problem}
\label{sec:bf-linear-kalman-hides-structural-split}

For a linear Gaussian state-space model,
\begin{align}
  x_t &= c + F x_{t-1} + u_t,\qquad u_t \sim \calN(0,Q),\\
  y_t &= a + H x_t + e_t,\qquad e_t \sim \calN(0,R),
\end{align}
the Kalman filter only needs the conditional mean and covariance of
$x_t\mid x_{t-1}$.  If a DSGE solution implies a singular or nearly singular
$Q$, the exact linear Gaussian likelihood remains well defined when the
collapsed representation preserves the same conditional law induced by the
structural transition.  Numerically stable implementations may then handle that
degenerate conditional Gaussian law through diffuse, square-root, solve-form,
or carefully regularized recursions as appropriate \citep{durbin2012time}.

This is why the following collapsed first-order representation can be
acceptable in an exact Kalman path:
\[
  Q = \eta \eta^\top,\qquad
  P_{t|t-1} = h_x P_{t-1|t-1} h_x^\top + Q.
\]
The key point is not that structural determinism disappears.  The key point is
that the collapsed linear representation still encodes the same one-step
conditional Gaussian law.  Deterministic-completion coordinates are not wrong
merely because they are included in $x_t$.  They obtain randomness through
their dependence on lagged stochastic coordinates, and their direct innovation
variance may be zero.  The exact linear Kalman recursion uses only the correct
first two conditional moments.

The danger is that this success can train the implementation culture into a
bad habit: it makes the state vector look like a homogeneous Gaussian vector.
That habit becomes wrong when the backend changes from exact linear filtering
to nonlinear sigma-point or particle filtering.  The relevant design decision
is then the integration space recorded in the filter metadata, not merely the
numerical dimension of $x_t$.

A related misunderstanding should be removed explicitly.  A deterministic-
completion coordinate need not have zero one-step predictive variance given the
lagged state.  It can inherit randomness through the current stochastic block.
In the Chapter~\ref{ch:bf-state-space-contracts} AR(2) lag-stack example, the
lag coordinate has no new innovation of its own but still remains random under
the lagged filtering distribution.  Likewise, in DSGE models with
\eqref{eq:bf-m-transition}--\eqref{eq:bf-k-transition}, $k_t$ may satisfy
\[
  \Var(k_t\mid x_{t-1}) > 0
\]
even though $k_t$ has no independent innovation once $(x_{t-1},\varepsilon_t)$
is fixed.

\section{Why Nonlinear Filters Must Respect the Structural Map}
\label{sec:bf-nonlinear-filters-need-structural-map}

Sigma-point and particle filters approximate a transition law, not just a
matrix covariance update.  In a generic nonlinear state-space model,
\[
  x_t = f_\theta(x_{t-1}, \varepsilon_t),\qquad
  y_t = g_\theta(x_t, e_t),
\]
the filter has to decide which variables are integrated over and which are
deterministic transforms.  In BayesFilter language, this is the choice of
integration space recorded in the filter metadata.  For DSGE models with
\eqref{eq:bf-m-transition}--\eqref{eq:bf-k-transition}, the natural nonlinear
prediction step is:
\begin{enumerate}[label=(\roman*)]
  \item represent the filtering distribution of the lagged structural state;
  \item draw, integrate, or choose quadrature points for the exogenous
    innovation $\varepsilon_t$ or current exogenous state $m_t$;
  \item compute each corresponding endogenous state $k_t$ using
    $T_k(k_{t-1},m_{t-1},m_t;\theta)$;
  \item evaluate the observation map on the resulting structurally valid
    points.
\end{enumerate}
The endogenous coordinates should not receive arbitrary independent
``sigma-point noise'' merely because they are entries of $x_t$.  If a sigma
point moves an accounting identity or predetermined endogenous state outside
the model-implied support of the structural transition, the filter is no longer
approximating the intended DSGE state-space model.  It is approximating a
different artificial model.  This is exactly the approximation boundary from
Chapter~\ref{ch:bf-sigma-point-filters}: quadrature may approximate the target
law, but injecting additional stochastic coordinates changes the target unless
that change is declared and justified.

Particle-filter literature makes the same distinction operationally.  A
particle filter propagates particles through the state transition and weights
them by the observation density \citep{gordon1993novel,doucet2001sequential,
andrieu2010particle}.  Written in the structural notation, a basic propagation
and weighting step is
\[
  x_t^{(i)} = F_\theta\!\left(x_{t-1}^{(a_i)},\varepsilon_t^{(i)}\right),
  \qquad
  w_t^{(i)} \propto p_\theta\!\left(y_t\mid x_t^{(i)}\right).
\]
When part of the transition is deterministic, particles should be propagated by
sampling the declared stochastic block and then deterministically completing the
remaining coordinates.  Adding artificial noise to deterministic coordinates can
be used only as a declared proposal or regularization device, with the
proposal/target correction and approximation status stated explicitly.  It is
not the default filtering law.

\section[Degenerate Transitions]{Degenerate Transitions Are Structural}
\label{sec:bf-degenerate-transitions-structural}

Degenerate transition covariance is not an exceptional corner case in DSGE
models.  It is often the mathematical expression of timing, identities, and
endogenous accumulation.  Capital, debt, lagged output growth, risk-premium
terms, and other predetermined variables may be deterministic conditional on
the previous state and current shocks.  A full covariance matrix with small positive diagonal padding can keep linear
algebra routines alive, but it does not by itself define the intended economic
transition.  BayesFilter therefore separates three ideas that are often
confused in practice:
\begin{enumerate}[label=(\roman*)]
  \item an exact degenerate transition law that already matches the structural
    model;
  \item numerical regularization used to evaluate that law stably; and
  \item an altered transition law that injects new stochastic degrees of
    freedom into deterministic-completion coordinates.
\end{enumerate}
Only the first two preserve the original structural model automatically.

The implementation policy is therefore:
\begin{itemize}
  \item distinguish structural determinism from numerical regularization;
  \item never let a nugget or PSD projection change the declared stochastic
    dimension without a written approximation label;
  \item expose deterministic blocks in diagnostics, rather than burying them
    under a full-rank covariance;
  \item test degenerate and rank-deficient transition cases directly.
\end{itemize}

This point is separate from the SVD spectral-gradient issue in
Chapter~\ref{ch:bf-svd-sigma-point}.  SVD factors may improve value-side
robustness for nearly singular covariances.  They do not repair a filter that
is propagating the wrong state-space law.

\section{Worked Structural UKF Example}
\label{sec:bf-structural-ukf-worked-example}

This section works through a one-step unscented update in a model with one
shock-driven state and one deterministic endogenous state.  The example is
small enough to expose the algebra, but it contains the same structural issue
as a DSGE model with exogenous shocks and endogenous predetermined variables.
The point of the example is not to prove a universal theorem about every UKF
variant.  The point is to show the BayesFilter structural-design rule for mixed
structural models: generate sigma points in the declared pre-transition
uncertainty variables, then complete deterministic coordinates pointwise.

Let
\begin{align}
  m_t &= \rho m_{t-1}+\sigma\varepsilon_t,
    & \varepsilon_t &\sim \calN(0,1), \label{eq:bf-ukf-example-m}\\
  k_t &= \phi k_{t-1}+\gamma m_t^2, \label{eq:bf-ukf-example-k}\\
  y_t &= m_t+k_t+e_t,
    & e_t &\sim \calN(0,R). \label{eq:bf-ukf-example-y}
\end{align}
The state is $x_t=(m_t,k_t)$, but only the innovation
$\varepsilon_t$ is new stochastic input at time $t$.  Conditional on
$(m_{t-1},k_{t-1},\varepsilon_t)$, the coordinate $k_t$ is deterministic.

Suppose the previous filtering distribution is Gaussian:
\[
  x_{t-1}\mid y_{1:t-1}
  \approx \calN(\mu_{t-1},P_{t-1}),
  \qquad
  \mu_{t-1}=
  \begin{bmatrix}\bar m\\ \bar k\end{bmatrix}.
\]
The structural UKF should augment the previous state with the innovation.  This
places sigma points in the same pre-transition uncertainty variables that define
the predictive law, rather than in an artificially padded post-transition full
state:
\[
  a_t =
  \begin{bmatrix}
    m_{t-1}\\ k_{t-1}\\ \varepsilon_t
  \end{bmatrix},
  \qquad
  a_t\mid y_{1:t-1}
  \approx
  \calN\left(
    \begin{bmatrix}\bar m\\ \bar k\\ 0\end{bmatrix},
    \begin{bmatrix}
      P_{t-1} & 0\\
      0 & 1
    \end{bmatrix}
  \right).
\]
Let $\{a_t^{(j)},w_j^{(m)},w_j^{(c)}\}_{j=0}^{2n_a}$ be the usual unscented
points and weights for this $n_a=3$ dimensional augmented Gaussian
\citep{julier1997new}.  For every point
$a_t^{(j)}=(m_{t-1}^{(j)},k_{t-1}^{(j)},\varepsilon_t^{(j)})$, the structural
transition is evaluated as
\begin{align}
  m_t^{(j)}
    &= \rho m_{t-1}^{(j)}+\sigma\varepsilon_t^{(j)},\\
  k_t^{(j)}
    &= \phi k_{t-1}^{(j)}+\gamma\left(m_t^{(j)}\right)^2,\\
  x_t^{(j)}
    &= \begin{bmatrix}m_t^{(j)}\\ k_t^{(j)}\end{bmatrix}.
\end{align}
This is the crucial step: the filter never creates an independent sigma-point
perturbation for $k_t$.  It creates sigma points for the previous state and
the current innovation, then computes $k_t$ by the structural map
\eqref{eq:bf-ukf-example-k}.

The predicted state moments are
\begin{align}
  \hat x_{t|t-1}
    &= \sum_j w_j^{(m)} x_t^{(j)},\\
  P_{xx,t}
    &= \sum_j w_j^{(c)}
       \left(x_t^{(j)}-\hat x_{t|t-1}\right)
       \left(x_t^{(j)}-\hat x_{t|t-1}\right)^\top .
\end{align}
For the observation equation \eqref{eq:bf-ukf-example-y}, define
\[
  z_t^{(j)}=m_t^{(j)}+k_t^{(j)}.
\]
Then the predicted observation, innovation variance, and cross covariance are
\begin{align}
  \hat z_t
    &= \sum_j w_j^{(m)} z_t^{(j)},\\
  S_t
    &= \sum_j w_j^{(c)}(z_t^{(j)}-\hat z_t)^2 + R,\\
  P_{xz,t}
    &= \sum_j w_j^{(c)}
       \left(x_t^{(j)}-\hat x_{t|t-1}\right)(z_t^{(j)}-\hat z_t).
\end{align}
The Gaussian measurement update is
\begin{align}
  v_t &= y_t-\hat z_t,\\
  K_t &= P_{xz,t}S_t^{-1},\\
  \hat x_{t|t} &= \hat x_{t|t-1}+K_tv_t,\\
  P_{t|t} &= P_{xx,t}-K_tS_tK_t^\top,
\end{align}
and the one-step approximate log likelihood contribution is
\[
  \ell_t
  =
  -\frac{1}{2}\left[
    \log(2\pi S_t)+\frac{v_t^2}{S_t}
  \right].
\]

\subsection{Numerical Illustration}
\label{sec:bf-structural-ukf-numeric}

Use the parameters
\[
  \rho=0.8,\quad \sigma=0.5,\quad \phi=0.7,\quad
  \gamma=0.4,\quad R=0.25,
\]
and previous filtering moments
\[
  \mu_{t-1}=(0,0)^\top,
  \qquad
  P_{t-1}=\operatorname{diag}(0.04,0.09).
\]
With the scaled unscented choice $\lambda=0$, the mean weight of the central
point is zero and the six off-center points each have mean weight $1/6$.
The augmented covariance is
\[
  \operatorname{diag}(0.04,0.09,1),
\]
so $\sqrt{(n_a+\lambda)P_a}$ has diagonal entries
$0.34641$, $0.51962$, and $1.73205$.  The seven sigma points and their
deterministic structural completion are:
\[
\begin{array}{c|rrr|rrr}
  j
  & m_{t-1}^{(j)} & k_{t-1}^{(j)} & \varepsilon_t^{(j)}
  & m_t^{(j)} & k_t^{(j)} & z_t^{(j)}\\
  \hline
  0 & 0       & 0        & 0        & 0        & 0        & 0\\
  1 & 0.34641& 0        & 0        & 0.27713 & 0.03072 & 0.30785\\
  2 & 0      & 0.51962  & 0        & 0       & 0.36373 & 0.36373\\
  3 & 0      & 0        & 1.73205  & 0.86603 & 0.30000 & 1.16603\\
  4 & -0.34641& 0       & 0        & -0.27713& 0.03072 & -0.24641\\
  5 & 0      & -0.51962 & 0        & 0       & -0.36373& -0.36373\\
  6 & 0      & 0        & -1.73205 & -0.86603& 0.30000 & -0.56603
\end{array}
\]
For example, the third positive innovation point gives
$m_t=0.5(1.73205)=0.86603$ and
$k_t=0.4(0.86603)^2=0.30000$.  No row contains an independently sampled
``new'' shock for $k_t$.

Using covariance weight $w_0^{(c)}=2$ for the central point, the structural
propagation gives
\[
  \hat x_{t|t-1}
  =
  \begin{bmatrix}
    0\\ 0.11024
  \end{bmatrix},
  \qquad
  P_{xx,t}
  =
  \begin{bmatrix}
    0.27560 & 0\\
    0 & 0.08657
  \end{bmatrix}.
\]
The predicted observation moments are
\[
  \hat z_t=0.11024,\qquad S_t=0.61217,\qquad
  P_{xz,t}=
  \begin{bmatrix}
    0.27560\\ 0.08657
  \end{bmatrix}.
\]
For an observation $y_t=0.30$,
\[
  v_t=0.18976,\qquad
  K_t=
  \begin{bmatrix}
    0.45020\\ 0.14141
  \end{bmatrix},
\]
and therefore
\[
  \hat x_{t|t}
  =
  \begin{bmatrix}
    0.08543\\ 0.13707
  \end{bmatrix},
  \qquad
  P_{t|t}
  =
  \begin{bmatrix}
    0.15152 & -0.03897\\
    -0.03897 & 0.07433
  \end{bmatrix}.
\]
The one-step log likelihood contribution is approximately
\[
  \ell_t=-0.70297.
\]

Now compare this with an off-manifold full-state approximation that adds an
independent artificial innovation variance $0.04$ directly to $k_t$ after the
structural map.  Nothing in
\eqref{eq:bf-ukf-example-m}--\eqref{eq:bf-ukf-example-y} justifies that extra
shock.  It changes the innovation variance from $S_t=0.61217$ to
$S_t=0.65217$ and changes the same observation's log likelihood contribution
from $-0.70297$ to $-0.73282$.  The numerical difference is small in this toy
calibration, but the conceptual difference is large: the latter filter is
approximating a different state-space model.

The lesson is not that UKF is wrong.  The lesson is that the UKF must be
applied to the correct structural transition.  Sigma points belong in the
declared stochastic integration space.  Endogenous deterministic coordinates
are completed by the model point by point.  In BayesFilter terms, this is a
design rule for mixed structural models and an approximation-label boundary for
any backend that instead enlarges the stochastic space.

\section{Pruned DSGE Filtering: A Minimal Correct Pattern}
\label{sec:bf-pruned-dsge-correct-pattern}

For pruned second-order DSGE approximations, the source DSGE project already
uses the standard first-order/second-order decomposition:
\begin{align}
  x_t^f &= h_x x_{t-1}^f + \eta\varepsilon_t,\\
  x_t^s &= h_x x_{t-1}^s
    + \frac{1}{2}H_{xx}(x_{t-1}^f\otimes x_{t-1}^f)
    + \frac{1}{2}h_{ss},\\
  y_t &= g_x(x_t^f+x_t^s)
    + \frac{1}{2}G_{xx}(x_t^f\otimes x_t^f)
    + \frac{1}{2}g_{ss} + e_t .
\end{align}
That decomposition is necessary, but it is not sufficient for large DSGE
filtering.  It separates perturbation order; it does not automatically
separate exogenous shock states from endogenous deterministic states inside
$x_t^f$.  A robust BayesFilter DSGE adapter should carry both decompositions:
\[
  \text{perturbation order: } (x^f,x^s),
  \qquad
  \text{structural timing: } (m,k).
\]

The implementation should therefore expose a structural transition interface
with at least:
\begin{itemize}
  \item exogenous-state indices and shock dimensions;
  \item endogenous/predetermined deterministic-state indices;
  \item the exogenous transition $T_m$;
  \item the deterministic endogenous transition $T_k$;
  \item the observation map and measurement-error law;
  \item a declaration of whether the backend integrates over shocks,
    exogenous states, or a full latent state approximation.
\end{itemize}
For CKF/UKF/SVD sigma-point filters, sigma points should be generated in the
declared stochastic integration space and then lifted through the structural
transition.  For particle filters, proposals should sample the stochastic
blocks and deterministically complete the endogenous block unless a different
proposal is explicitly specified.

\section{Current Source-Project Failure Mode}
\label{sec:bf-current-source-project-failure-mode}

The recent source-project audit found a concrete instance of this design
failure.  The default DSGE SVD sigma-point adapter uses a collapsed transition
of the form
\[
  Q_w = \eta\eta^\top + \epsilon I,\qquad
  x_{t|t-1}^f = h_x x_{t-1|t-1}^f,
\]
and propagates the covariance as
\[
  P_{t|t-1} = h_x P_{t-1|t-1} h_x^\top + Q_w.
\]
It tracks the pruned $x^s$ correction separately, but it does not expose a
partition between exogenous and endogenous coordinates inside $x^f$.  For the smallest NK example,
where the states are only the exogenous demand and monetary-policy shocks,
this is not a serious modeling error.  For Rotemberg NK, SGU, EZ, and
NAWM-scale DSGE models, the state vector contains deterministic or
endogenous/predetermined components.  Treating those models as if every state
coordinate were an exchangeable noisy Gaussian coordinate is a structural
filtering bug, not a mere numerical tuning problem.

\section{Required Tests Before Any DSGE Nonlinear Filter Claim}
\label{sec:bf-structural-deterministic-required-tests}

No future BayesFilter report should claim that a nonlinear DSGE filter is
value-valid, gradient-valid, sampler-usable, converged, or production-ready in
the sense of Chapter~\ref{ch:bf-production-checklist} unless the following
structural tests are passed or explicitly waived with written justification:
\begin{enumerate}
  \item \textbf{Metadata test:} the model exposes exogenous, endogenous, and
    deterministic state blocks, and their dimensions match the shock-impact
    matrix and transition maps.
  \item \textbf{Constraint-support test:} propagated sigma points or particles
    satisfy deterministic identities to numerical tolerance.
  \item \textbf{Linear recovery test:} when the structural transition is
    linear, the structural nonlinear filter recovers the exact Kalman
    likelihood up to the declared approximation tolerance.
  \item \textbf{Degenerate-transition test:} zero-innovation endogenous rows
    remain deterministic except for explicitly declared numerical
    regularization.
  \item \textbf{Toy DSGE test:} a controlled model with
    \[
      m_t = \rho m_{t-1}+\sigma\varepsilon_t,\qquad
      k_t=f(k_{t-1},m_{t-1},m_t),
    \]
    compares structural propagation against a dense quadrature or
    high-particle reference.
  \item \textbf{Case-study test:} Rotemberg NK, SGU, and any EZ/NAWM target
    must document which state coordinates are shock-driven, which are
    deterministic conditional on the shock path, and how the filter honors
    that timing.
\end{enumerate}

\section{Common Misunderstandings}
\label{sec:bf-structural-common-misunderstandings}

The most common reader errors in this area are:
\begin{enumerate}[label=(\roman*)]
  \item a deterministic endogenous state need not have zero one-step predictive
    variance; it can inherit randomness through the declared stochastic block;
  \item a singular or rank-deficient transition covariance is not, by itself, a
    modeling error in the exact linear-Gaussian case;
  \item numerical regularization is not the same thing as declaring a new
    stochastic state coordinate;
  \item a full-state nonlinear approximation is not automatically forbidden, but
    it must be labeled as an approximation if it enlarges the stochastic space
    beyond the structural transition contract.
\end{enumerate}


The rule is simple:
\begin{center}
\fbox{\parbox{0.88\linewidth}{
For nonlinear DSGE filtering, the stochastic integration variables are not
automatically the entries of the state vector.  They are the variables declared
stochastic by the structural transition.  Deterministic endogenous states must
be computed, not noised.
}}
\end{center}

If an implementation violates this rule, it must be labeled as an artificial
approximation and justified against a reference.  Otherwise, the filter is
solving the wrong state-space model.  This policy should be treated as a release gate for any BayesFilter DSGE
backend intended for HMC, industrial forecasting, or peer-reviewed empirical
work.  Failing one of these tests does not automatically make a backend useless,
but it does mean the strongest claim labels from
Chapter~\ref{ch:bf-production-checklist} are not yet justified.
